\chapter{Customization}
\label{ch:customize}

\index{customization}

Acme deliberately has very little to configure.  There are no
configuration files, no themes, no plugin system.  This is by design:
the power comes from composing external tools, not from configuring the
editor.  Nevertheless, there are several things you can adjust.

\section{Fonts}
\label{sec:fonts}

\index{fonts}

Acme maintains two fonts: a proportional (variable-width) font and a
fixed-width font.  You can set both at startup and switch between them
per window.

\subsection{Startup Fonts}

\begin{lstlisting}[language=bash]
acme -f 'DejaVuSans/14a/font' -F 'DejaVuSansMono/14a/font'
\end{lstlisting}

The font name format is a fontconfig name followed by a size.  The
\texttt{a} suffix indicates anti-aliased rendering.

\subsection{Per-Window Fonts}

Middle-click \cmd{Font} in any window tag to toggle between the
proportional and fixed-width font.

To set a specific font for one window, type:

\begin{Verbatim}[frame=single]
Font /mnt/font/GoMono/13a/font
\end{Verbatim}

\noindent
and middle-click the whole thing.

\subsection{Available Fonts}

Acme discovers fonts through fontconfig.  Any TrueType or OpenType font
installed on your system is available.  To list installed fonts:

\begin{lstlisting}[language=bash]
fc-list : family | sort -u
\end{lstlisting}

\section{Tab Width}
\label{sec:tabwidth}

\index{tab width}

The default tab width is 4 columns.  To change it:

\begin{description}[style=nextline]
\item[Per window:]
Middle-click \cmd{Tab 8} (or any number) in the window tag.

\item[At startup:]
Set the \env{tabstop} environment variable:
\begin{lstlisting}[language=bash]
tabstop=8 acme
\end{lstlisting}
\end{description}

\section{Auto-Indentation}
\label{sec:autoindent}

\index{auto-indent}

Auto-indent is off by default.  When enabled, pressing \key{Enter}
copies the leading whitespace of the current line onto the new line.

\begin{description}[style=nextline]
\item[At startup:]
Use the \cmd{-a} flag:
\begin{lstlisting}[language=bash]
acme -a
\end{lstlisting}

\item[Per window:]
Middle-click \cmd{Indent on} or \cmd{Indent off} in the tag.

\item[Globally:]
Middle-click \cmd{Indent ON} or \cmd{Indent OFF} (uppercase).
\end{description}

\section{Plumbing Rules}
\label{sec:plumbrules}

\index{plumbing!customization}

The most significant customization you can make is writing custom
plumbing rules.  Create \file{\textasciitilde/lib/plumbing} with your
rules.  See Chapter~\ref{ch:plumbing} for the full syntax and examples.

Common customizations:

\begin{itemize}
\item Add rules for your programming language's error format.
\item Add rules for documentation URLs (e.g., Go import paths,
      Rust crate docs).
\item Add rules for your project's ticket system (e.g., clicking
      ``JIRA-1234'' opens the ticket in a browser).
\item Change the file opener command.
\end{itemize}

\section{Shell}
\label{sec:shellconfig}

\index{shell configuration}

\subsection{Command Shell}

The shell used for executing commands (middle-click on text) is
controlled by \env{acmeshell}:

\begin{lstlisting}[language=bash]
acmeshell=bash acme
\end{lstlisting}

If not set, acme defaults to \cmd{sh}.

\subsection{Win Shell}

The \cmd{Win} command uses \env{SHELL}:

\begin{lstlisting}[language=bash]
SHELL=/bin/zsh acme
\end{lstlisting}

\section{Helper Scripts}
\label{sec:helpers}

\index{helper scripts}

Since acme sets \env{winid} and \env{\%} in the environment of child
processes, you can write scripts that interact with the current window.

\subsection{Example: Format on Save}

Add to a tag:

\begin{Verbatim}[frame=single]
Fmt
\end{Verbatim}

\noindent
Where \cmd{Fmt} is a script in your PATH:

\begin{lstlisting}[language=bash]
#!/bin/sh
# Fmt: format the current file and reload
case "$%" in
    *.go) gofmt -w "$%" ;;
    *.py) black "$%" ;;
    *.c)  indent -kr "$%" ;;
    *.rs) rustfmt "$%" ;;
esac
\end{lstlisting}

After running \cmd{Fmt}, middle-click \cmd{Get} to reload the
formatted file.

\subsection{Example: Run Tests for Current File}

\begin{lstlisting}[language=bash]
#!/bin/sh
# Test: run tests related to the current file
case "$%" in
    *_test.go) go test -run . -v "$(dirname "$%")" ;;
    *.py)      python -m pytest "$%" -v ;;
    *.c)       make test ;;
esac
\end{lstlisting}

\subsection{Example: Open Documentation}

\begin{lstlisting}[language=bash]
#!/bin/sh
# Doc: open documentation for the word argument
word="$1"
case "$%" in
    *.go) xdg-open "https://pkg.go.dev/search?q=$word" ;;
    *.py) xdg-open "https://docs.python.org/3/search.html?q=$word" ;;
    *.rs) xdg-open "https://docs.rs/releases/search?query=$word" ;;
esac
\end{lstlisting}

\section{Window Size}
\label{sec:winsize}

\index{window size}

Set the initial window size with the \cmd{-W} flag:

\begin{lstlisting}[language=bash]
acme -W 1400x900
\end{lstlisting}

\section{Scrollbar Direction}
\label{sec:scrolldir}

\index{scrollbar!direction}

The \cmd{-r} flag reverses the scrollbar button mapping.  By default,
\button{1} in the scrollbar scrolls up and \button{3} scrolls down.
With \cmd{-r}, this is reversed.

\begin{lstlisting}[language=bash]
acme -r
\end{lstlisting}

\section{Startup Script}

A common pattern is to create a shell alias or script that starts acme
with your preferred settings:

\begin{lstlisting}[language=bash]
#!/bin/sh
# ~/bin/a - start acme with my preferences
export tabstop=4
export acmeshell=bash
exec acme -a \
    -f 'GoRegular/13a/font' \
    -F 'GoMono/13a/font' \
    "$@"
\end{lstlisting}

There is no \file{.acmerc} or configuration file by design.  A shell
script gives you the same effect with full transparency.
